% Unit 058 English translation of chapter4.tex lines 1681--1840.
% Source: Methods of Algebra, Volume 2, commit
% 9a5803ff77dd3257484cb177f851a73770a59dd3 (CC BY 4.0).
% stable-unit-id: o014.aljabr2.chapter4.bounded-derived-functors
% Coverage: complete Section 4.8; stops before Section 4.9 at line 1841.

% segment-id: o014.li2.chapter4.bounded-derived-functors.h001
\section{Bounded Derived Functors}\label{sec:bounded-derived-functor}
\hypertarget{o014.aljabr2.chapter4.bounded-derived-functors}{ }

% segment-id: o014.li2.chapter4.bounded-derived-functors.p001
This section explains how to study ${}^+\mathrm{R}F$ and
${}^-\mathrm{L}F$ by means of injective and projective resolutions. The case
of unbounded derived categories is deferred to \S\ref{sec:unbounded-derived}.

% segment-id: o014.li2.chapter4.bounded-derived-functors.c001
\begin{convention}
	This section focuses on $\star\in\{+,-\}$ and adopts the following
	conventions.
	\begin{itemize}
		\item For right derived functors we always take $\star=+$ and use the
		abbreviation
		$\mathrm{R}F:={}^+\mathrm{R}F:
		\cate{D}^+(\mathcal{A})\to\cate{D}^+(\mathcal{A}')$.
		\item For left derived functors we always take $\star=-$ and use the
		abbreviation
		$\mathrm{L}F:={}^-\mathrm{L}F:
		\cate{D}^-(\mathcal{A})\to\cate{D}^-(\mathcal{A}')$.
	\end{itemize}
	In view of Lemma~\ref{prop:derived-functor-res}, these conventions cause no
	confusion.
\end{convention}

% segment-id: o014.li2.chapter4.bounded-derived-functors.dp001
\begin{definition-proposition}\label{prop:F-injective-criterion}
	\index{subcategory of type I}
	\index{subcategory of type P}
	Let $\mathcal{A}^\flat$ be a full additive subcategory of $\mathcal{A}$.
	We call $\mathcal{A}^\flat$ \emph{of type I} relative to $F$ if the
	following conditions hold; in this case
	$\cate{K}^+(\mathcal{A}^\flat)$ is an $F$-injective subcategory of
	$\cate{K}^+(\mathcal{A})$:
	\begin{enumerate}[({I}1)]
		\item for every $X\in\Obj(\mathcal{A})$, there is a monomorphism
		$X\hookrightarrow I$ with $I\in\Obj(\mathcal{A}^\flat)$;
		\item if $0\to X'\to X\to X''\to0$ is a short exact sequence in
		$\mathcal{A}$ and $X',X\in\Obj(\mathcal{A}^\flat)$, then
		$X''\in\Obj(\mathcal{A}^\flat)$ and
		$0\to FX'\to FX\to FX''\to0$ is exact.
	\end{enumerate}
	Dually, we call $\mathcal{A}^\flat$ \emph{of type P} relative to $F$ if
	the following conditions hold; in this case
	$\cate{K}^-(\mathcal{A}^\flat)$ is an $F$-projective subcategory of
	$\cate{K}^-(\mathcal{A})$:
	\begin{enumerate}[({P}1)]
		\item for every $X\in\Obj(\mathcal{A})$, there is an epimorphism
		$P\twoheadrightarrow X$ with $P\in\Obj(\mathcal{A}^\flat)$;
		\item if $0\to X'\to X\to X''\to0$ is a short exact sequence in
		$\mathcal{A}$ and $X,X''\in\Obj(\mathcal{A}^\flat)$, then
		$X'\in\Obj(\mathcal{A}^\flat)$ and
		$0\to FX'\to FX\to FX''\to0$ is exact.
	\end{enumerate}
\end{definition-proposition}
% segment-id: o014.li2.chapter4.bounded-derived-functors.pf001
\begin{proof}
	It suffices to treat the type-I version. First,
	$\cate{K}^+(\mathcal{A}^\flat)$ is a triangulated subcategory of
	$\cate{K}^+(\mathcal{A})$ (Corollary~\ref{prop:subcat-K-triangulated}).
	Theorem~\ref{prop:resolution-existence} shows that for every
	$X\in\Obj(\cate{K}^+(\mathcal{A}))$ there is a quasi-isomorphism $X\to I$
	with $I\in\cate{K}^+(\mathcal{A}^\flat)$. This is the resolution condition
	for an $F$-injective subcategory. It remains to prove that if
	$X\in\cate{K}^+(\mathcal{A}^\flat)$ is acyclic, then $FX$ is also acyclic.

	For every $n\in\Z$, there is a short exact sequence
	$0\to\Ker(d^n)\to X^n\xrightarrow{d^n}\Ker(d^{n+1})\to0$. If $n\ll0$,
	then $\Ker(d^n)=X^n=0$. Applying (I2) recursively for all $n$ gives
	$\Ker(d^n)\in\Obj(\mathcal{A}^\flat)$ together with short exact sequences
	in $\mathcal{A}'$
	\[ 0 \to F\Ker(d^n) \to FX^n \xrightarrow{Fd^n} F\Ker(d^{n+1}) \to 0. \]
	Splicing these short exact sequences shows that $FX$ is acyclic.
\end{proof}

% segment-id: o014.li2.chapter4.bounded-derived-functors.p002
The terms type I and type P are merely convenient labels. There is a complete
characterization of the subcategories $\mathcal{A}^\flat$ that give
$F$-injective or $F$-projective subcategories; see
\cite[Proposition 13.3.5]{KS06}.

% segment-id: o014.li2.chapter4.bounded-derived-functors.p003
The next result shows that derived functors can be computed by taking
resolutions in a full subcategory of type I or P.

% segment-id: o014.li2.chapter4.bounded-derived-functors.t001
\begin{theorem}\label{prop:derived-via-resolution}
	Suppose that $\mathcal{A}$ has a full additive subcategory
	$\mathcal{A}^\flat$ of type I (respectively, type P) relative to the
	additive functor $F:\mathcal{A}\to\mathcal{A}'$. Then the derived functor
	$\mathrm{R}F$ (respectively, $\mathrm{L}F$) exists. More precisely:
	\begin{enumerate}[(i)]
		\item for every $X\in\Obj(\cate{K}^+(\mathcal{A}))$ (respectively,
		$X\in\Obj(\cate{K}^-(\mathcal{A}))$), there is a quasi-isomorphism
		$X\to I$ (respectively, $P\to X$) with
		$I\in\Obj(\cate{K}^+(\mathcal{A}^\flat))$ (respectively,
		$P\in\Obj(\cate{K}^-(\mathcal{A}^\flat))$);
		\item for the quasi-isomorphism in (i), the morphism in
		\eqref{eqn:derived-can-morphism} gives
		\[ \mathrm{R}F(QX) \rightiso \mathrm{R}F(QI) \leftiso Q' \cate{K}^+ F(I), \quad \mathrm{L}F(QX) \leftiso \mathrm{L}F(QP) \rightiso Q' \cate{K}^- F(P); \]
		\item $\mathrm{R}F$ restricts to
		$\cate{D}^{\geq0}(\mathcal{A})\to\cate{D}^{\geq0}(\mathcal{A}')$
		(respectively, $\mathrm{L}F$ restricts to
		$\cate{D}^{\leq0}(\mathcal{A})\to\cate{D}^{\leq0}(\mathcal{A}')$);
		\item suppose that $F$ is left exact (respectively, right exact). If
		$X\in\Obj(\mathcal{A})$, the morphism in
		\eqref{eqn:derived-can-morphism} gives an isomorphism
		$FX\rightiso\mathrm{R}^0F(QX)$ (respectively,
		$FX\leftiso\mathrm{L}_0F(QX)$);
		\item if $X\in\Obj(\mathcal{A}^\flat)$, then, for $n>0$,
		$\mathrm{R}^nF(X)=0$ (respectively, $\mathrm{L}_nF(X)=0$).
	\end{enumerate}
\end{theorem}
% segment-id: o014.li2.chapter4.bounded-derived-functors.pf002
\begin{proof}
	It suffices to consider $\mathrm{R}F$. Since
	$\cate{K}^+(\mathcal{A}^\flat)$ is $F$-injective, a quasi-isomorphism
	$X\to I$ always exists. The existence of $\mathrm{R}F$ and the isomorphism
	$\mathrm{R}F(QX)\simeq Q'\cate{K}^+F(I)$ merely restate
	Proposition~\ref{prop:localization-injective}. This proves (i) and (ii).

	For (iii), if $X\in\Obj(\cate{D}^{\geq0}(\mathcal{A}))$, we may assume
	that $X$ comes from an object of $\cate{C}^{\geq0}(\mathcal{A})$.
	Theorem~\ref{prop:resolution-existence} allows us to choose a
	quasi-isomorphism $X\to I$ with
	$I\in\Obj(\cate{C}^{\geq0}(\mathcal{A}^\flat))$. Hence
	\[ \mathrm{R}F(QX) \simeq Q' \cate{K}^+ F(I) \in \Obj\left( \cate{D}^{\geq 0}(\mathcal{A}')\right). \]
	If, in addition, $X\in\Obj(\mathcal{A}^\flat)$, the quasi-isomorphism
	$X\to I$ may be taken to be $\identity_X$. In that case
	$\mathrm{R}F(QX)=Q'\cate{K}^+F(X)=Q'FX$ is concentrated in degree zero.
	This also proves (v).

	Finally, consider (iv). Suppose that $F$ is left exact. For
	$X\in\Obj(\mathcal{A})$, choose an exact sequence
	$0\to X\to I^0\to I^1\to\cdots$, with every
	$I^n\in\Obj(\mathcal{A}^\flat)$. Put $I=[I^0\to I^1\to\cdots]$.
	By (ii) and left exactness,
	$FX\rightiso\Ker[Fd_I^0:F(I^0)\to F(I^1)]
	\simeq\mathrm{R}^0F(X)$, as desired.
\end{proof}

% segment-id: o014.li2.chapter4.bounded-derived-functors.e001
\begin{example}[Injective and projective objects]\label{eg:injective-gives-F-injectives}
	Suppose that $\mathcal{A}$ has enough injective objects, and let
	$\mathcal{A}^\flat$ be the full additive subcategory of injective objects.
	It is of type I relative to every $F$: condition (I1) holds by assumption,
	while for (I2), Lemma~\ref{prop:inj-proj-split} shows that the short exact
	sequence in that condition splits. Thus $X'\oplus X''$ is injective, and
	Lemma~\ref{prop:prod-injective-projective} implies that $X''$ is injective
	as well.

	Dually, if $\mathcal{A}$ has enough projective objects, the full additive
	subcategory consisting of those objects is of type P relative to every $F$.
\end{example}

% segment-id: o014.li2.chapter4.bounded-derived-functors.p004
Example~\ref{eg:injective-gives-F-injectives} connects this discussion with
the classical definition of derived functors in \S\ref{sec:derived-primer}.
Suppose that $\mathcal{A}$ has enough injective objects. Consider the
restriction $\mathrm{R}^nF|_{\mathcal{A}}$ of the right derived functor. Since
a short exact sequence in $\mathcal{A}$ extends canonically to a distinguished
triangle (Proposition~\ref{prop:ses-vs-triangle}), the long exact sequence
immediately makes $(\mathrm{R}^nF|_{\mathcal{A}})_{n\geq0}$ a cohomological
$\delta$-functor in the sense of Definition~\ref{def:delta-functor}. Dually,
the left derived functors form the homological $\delta$-functor
$(\mathrm{L}_nF|_{\mathcal{A}})_{n\geq0}$.

% segment-id: o014.li2.chapter4.bounded-derived-functors.pr001
\begin{proposition}
	Suppose that $F$ is left exact (respectively, right exact) and that
	$\mathcal{A}$ has enough injective (respectively, projective) objects. Then
	$\mathrm{R}^nF:\cate{C}^+(\mathcal{A})\to\mathcal{A}'$ (respectively,
	$\mathrm{L}_nF:\cate{C}^-(\mathcal{A})\to\mathcal{A}'$) agrees with the
	version previously defined in Definition~\ref{def:derived-primer}.

	More generally, if there is a subcategory $\mathcal{A}^\flat$ of type I
	(respectively, type P), as in
	Definition--Proposition~\ref{prop:F-injective-criterion}, then
	$(\mathrm{R}^nF|_{\mathcal{A}})_{n\geq0}$ (respectively,
	$(\mathrm{L}_nF|_{\mathcal{A}})_{n\geq0}$) is a universal
	$\delta$-functor in the sense of
	Definition~\ref{def:universal-delta-functor}.
\end{proposition}
% segment-id: o014.li2.chapter4.bounded-derived-functors.pf003
\begin{proof}
	It suffices to treat the left-exact case. For
	$X\in\Obj(\cate{C}^+(\mathcal{A}))$, choose an injective resolution
	$X\to I$. Theorem~\ref{prop:derived-via-resolution} (ii) gives canonical
	isomorphisms
	$\mathrm{R}^nF(QX)\simeq\Hm^n\cate{K}^+F(I)$ for every $n\in\Z$. The
	right-hand side is precisely the version of the derived functor from
	Definition~\ref{def:derived-primer}.

	More generally, suppose that $\mathcal{A}^\flat$ is of type I relative to
	$F$. To prove the universality of
	$(\mathrm{R}^nF|_{\mathcal{A}})_{n\geq0}$, by
	Proposition~\ref{prop:erasable-univ-delta} it suffices to show that
	$\mathrm{R}^nF|_{\mathcal{A}}$ is effaceable for $n>0$. This follows
	directly from condition (I1) in
	Definition--Proposition~\ref{prop:F-injective-criterion} and
	Theorem~\ref{prop:derived-via-resolution} (v).
\end{proof}

% segment-id: o014.li2.chapter4.bounded-derived-functors.p005
To derive a composite of functors, one often needs a larger class of objects
for resolutions. We need the following concept.

% segment-id: o014.li2.chapter4.bounded-derived-functors.d001
\begin{definition}\label{def:F-acyclic-object}
	\index{$F$-acyclic object}
	Suppose that $\mathrm{R}F$ (respectively, $\mathrm{L}F$) exists. An object
	$X\in\Obj(\mathcal{A})$ is called \emph{$F$-acyclic} if
	$\mathrm{R}F(X)$ (respectively, $\mathrm{L}F(X)$) is an object of
	$\cate{D}^{\geq0}(\mathcal{A}')\cap\cate{D}^{\leq0}(\mathcal{A}')
	\simeq\mathcal{A}'$.
\end{definition}

% segment-id: o014.li2.chapter4.bounded-derived-functors.p006
The next result shows that derived functors can be computed using
$F$-acyclic resolutions; compare Corollary~\ref{prop:acyclic-resolution}.

% segment-id: o014.li2.chapter4.bounded-derived-functors.co001
\begin{corollary}\label{prop:derived-via-F-acyclic}
	Let $\mathcal{A}^\flat$ be a full additive subcategory of $\mathcal{A}$ of
	type I (respectively, type P) relative to $F$. Define
	$\mathcal{A}^\natural$ to be the full additive subcategory consisting of all
	$F$-acyclic objects of $\mathcal{A}$. Then:
	\begin{enumerate}[(i)]
		\item $\mathcal{A}^\natural\supset\mathcal{A}^\flat$;
		\item $\mathcal{A}^\natural$ is also of type I (respectively, type P);
		\item for every quasi-isomorphism $X\to I$ (respectively, $P\to X$),
		with $I\in\Obj(\cate{K}^+(\mathcal{A}^\natural))$ (respectively,
		$P\in\Obj(\cate{K}^-(\mathcal{A}^\natural))$), the morphism in
		\eqref{eqn:derived-can-morphism} gives
		\[ \mathrm{R}F(QX) \rightiso \mathrm{R}F(QI) \leftiso Q' \cate{K}^+ F(I), \quad \mathrm{L}F(QX) \leftiso \mathrm{L}F(QP) \rightiso Q' \cate{K}^- F(P). \]
	\end{enumerate}
\end{corollary}
% segment-id: o014.li2.chapter4.bounded-derived-functors.pf004
\begin{proof}
	Theorem~\ref{prop:derived-via-resolution} (v) implies (i), so condition (I1)
	or (P1) also holds for $\mathcal{A}^\natural$. For (I2), suppose that
	$0\to X'\to X\to X''\to0$ is a short exact sequence in $\mathcal{A}$ and
	that $X',X$ are $F$-acyclic. The long exact sequence from
	Theorem~\ref{prop:derived-long-exact} gives exact sequences
	\[ \underbracket{\mathrm{R}^n F(X)}_{= 0} \to \mathrm{R}^n F(X'') \to \underbracket{\mathrm{R}^{n+1}F(X')}_{= 0}, \quad n \in \Z_{\geq 1} , \]
	so $X''$ is $F$-acyclic. On the other hand, the $n=0$ portion of the long
	exact sequence, together with $\mathrm{R}^1F(X')=0$, shows that
	$0\to FX'\to FX\to FX''\to0$ is exact. The proof of (P2) is dual. This
	proves (ii).

	Finally, (iii) follows by applying
	Theorem~\ref{prop:derived-via-resolution} (ii) to
	$\mathcal{A}^\natural$.
\end{proof}

% segment-id: o014.li2.chapter4.bounded-derived-functors.p007
We now discuss deriving a composite of functors. The basic tool is
Theorem~\ref{prop:localization-triangulated-composite}.

% segment-id: o014.li2.chapter4.bounded-derived-functors.t002
\begin{theorem}[Deriving a composite of functors]\label{prop:derived-composite}
	Consider additive functors between abelian categories
	\[ \mathcal{A} \xrightarrow{F} \mathcal{A}' \xrightarrow{F'} \mathcal{A}''. \]
	Suppose that there are full subcategories $\mathcal{A}^\flat$ of
	$\mathcal{A}$ and $(\mathcal{A}')^\flat$ of $\mathcal{A}'$, of type I
	relative to $F$ and $F'$, respectively, in the sense of
	Definition--Proposition~\ref{prop:F-injective-criterion}. If $F$ sends the
	objects of $\mathcal{A}^\flat$ to $F'$-acyclic objects of $\mathcal{A}'$,
	then the canonical morphism
	$\mathrm{R}(F'F)\to(\mathrm{R}F')(\mathrm{R}F)$ is an isomorphism.

	There is a corresponding version for the left derived functors
	$\mathrm{L}F$, $\mathrm{L}F'$, and the canonical morphism
	$(\mathrm{L}F')(\mathrm{L}F)\to\mathrm{L}(F'F)$, using subcategories of
	type P.
\end{theorem}
% segment-id: o014.li2.chapter4.bounded-derived-functors.pf005
\begin{proof}
	It suffices to treat right derived functors. Let
	$(\mathcal{A}')^\natural$ be the full additive subcategory of $F'$-acyclic
	objects in $\mathcal{A}'$. Corollary~\ref{prop:derived-via-F-acyclic},
	together with
	Definition--Proposition~\ref{prop:F-injective-criterion}, shows that
	$\cate{K}^+(\mathcal{A}^\flat)$ and
	$\cate{K}^+((\mathcal{A}')^\natural)$ are, respectively, $F$-injective and
	$F'$-injective subcategories of $\cate{K}^+(\mathcal{A})$ and
	$\cate{K}^+(\mathcal{A}')$. Apply
	Theorem~\ref{prop:localization-triangulated-composite}.
\end{proof}

% segment-id: o014.li2.chapter4.bounded-derived-functors.p008
Using the amplitude from Definition~\ref{def:functor-amplitude}, we can define
the dimensions of the right and left derived functors of $F$.

% segment-id: o014.li2.chapter4.bounded-derived-functors.dp002
\begin{definition-proposition}\label{def:dim-functor}
	\index{dimension}
	Suppose that $F:\mathcal{A}\to\mathcal{A}'$ is left exact (respectively,
	right exact), and that the derived functor $\mathrm{R}F$ (respectively,
	$\mathrm{L}F$) exists and is determined as in
	Theorem~\ref{prop:derived-via-resolution}. There is a
	$d\in\Z_{\geq0}\sqcup\{+\infty\}$ such that the amplitude of the derived
	functor is contained in $[0,d]$ (respectively, $[-d,0]$). The infimum of
	all such $d$ is called the \emph{dimension} of $\mathrm{R}F$
	(respectively, $\mathrm{L}F$).
\end{definition-proposition}
% segment-id: o014.li2.chapter4.bounded-derived-functors.pf006
\begin{proof}
	It suffices to treat the left-exact case.
	Theorem~\ref{prop:derived-via-resolution} implies that $\mathrm{R}F$
	restricts to $\mathcal{A}\to\cate{D}^{\geq0}(\mathcal{A}')$, so its
	amplitude is contained in an interval $[0,d]$.
\end{proof}

% segment-id: o014.li2.chapter4.bounded-derived-functors.p009
Consequently, if $\mathrm{R}F$ (respectively, $\mathrm{L}F$) exists and has
finite dimension, it restricts to
$\cate{D}^{\bdd}(\mathcal{A})\to\cate{D}^{\bdd}(\mathcal{A}')$.

% segment-id: o014.li2.chapter4.bounded-derived-functors.e002
\begin{example}[$\Tor$-dimension]\label{eg:Tor-dimension}
	\index{dimension!$\operatorname{Tor}$}
	Let $R$ be a ring and $X$ a right $R$-module. Consider the right-exact
	functor from $R\dcate{Mod}$ to $\cate{Ab}$ that sends
	$Y\mapsto X\otimes Y:=X\dotimes{R}Y$. Since $R\dcate{Mod}$ has enough
	projective objects, $\mathrm{L}(X\otimes\cdot)$ exists, and
	$\mathrm{L}_n(X\otimes\cdot)$ is precisely
	$\Tor^R_n(X,\cdot)$ from
	Definition--Proposition~\ref{def:Tor-classical}. The dimension of
	$\mathrm{L}(X\otimes\cdot)$ is called the \emph{$\Tor$-dimension} of
	$X$. A similar definition applies to left $R$-modules, and
	Proposition~\ref{prop:Tor-symmetry} implies that the two agree when $R$ is
	commutative. The derived tensor-product functor is discussed further in
	\S\ref{sec:otimesL}.
\end{example}

% segment-id: o014.li2.chapter4.bounded-derived-functors.p010
Now fix abelian categories $\mathcal{A}_1$, $\mathcal{A}_2$, and
$\mathcal{B}$, and an additive bifunctor
\[ F: \mathcal{A}_1 \times \mathcal{A}_2 \to \mathcal{B}. \]
We do not require $\mathcal{B}$ to have countable coproducts or products,
because $\cate{C}^2F$ restricts to
\begin{gather*}
	\cate{C}^2 F: \cate{C}^{\pm}(\mathcal{A}_1) \times \cate{C}^{\pm} (\mathcal{A}_2) \to \cate{C}^2_f(\mathcal{B}).
\end{gather*}
On $\cate{C}^2_f(\mathcal{B})$, the total-complex functor
$\tot_{\oplus}=\tot_{\Pi}$ is always defined; denote this functor by
$\tot$. Likewise, denote $\cate{C}_{\oplus}F=\cate{C}_{\Pi}F$ by
$\cate{C}F$, and $\cate{K}_{\oplus}F=\cate{K}_{\Pi}F$ by $\cate{K}F$.

% segment-id: o014.li2.chapter4.bounded-derived-functors.p011
To study the derived bifunctors of Definition~\ref{def:derived-bifunctor}, we
consider triangulated subcategories $\mathcal{I}_i$ (respectively,
$\mathcal{P}_i$) of $\cate{K}^+(\mathcal{A}_i)$ (respectively,
$\cate{K}^-(\mathcal{A}_i)$). As usual, the sign is determined by whether we
are considering the right or left derived bifunctor, and $i=1,2$. How should
we choose the pair $(\mathcal{I}_1,\mathcal{I}_2)$ or
$(\mathcal{P}_1,\mathcal{P}_2)$?
Definition--Proposition~\ref{prop:F-injective-criterion} and the techniques of
\S\ref{sec:double-cplx-coh} give the following answer.

% segment-id: o014.li2.chapter4.bounded-derived-functors.l001
\begin{lemma}\label{prop:RF-bifunctor}
	Let $\mathcal{A}_i^\flat\subset\mathcal{A}_i$ be full additive
	subcategories, $i=1,2$. If the following conditions hold simultaneously,
	then the pair
	$\bigl(\cate{K}^+(\mathcal{A}_1^\flat),
	\cate{K}^+(\mathcal{A}_2^\flat)\bigr)$ (respectively,
	$\bigl(\cate{K}^-(\mathcal{A}_1^\flat),
	\cate{K}^-(\mathcal{A}_2^\flat)\bigr)$) is $F$-injective (respectively,
	$F$-projective) in the sense of
	Definition~\ref{def:F-injective-bifunctor}:
	\begin{itemize}
		\item for fixed $X_1\in\Obj(\mathcal{A}_1^\flat)$,
		$\mathcal{A}_2^\flat$ is of type I (respectively, type P) relative to
		$F(X_1,\cdot)$;
		\item for fixed $X_2\in\Obj(\mathcal{A}_2^\flat)$,
		$\mathcal{A}_1^\flat$ is of type I (respectively, type P) relative to
		$F(\cdot,X_2)$.
	\end{itemize}
	In particular, if $\mathcal{A}_1$ and $\mathcal{A}_2$ each have enough
	injective (respectively, projective) objects, then $\mathrm{R}F$
	(respectively, $\mathrm{L}F$) exists.
\end{lemma}
% segment-id: o014.li2.chapter4.bounded-derived-functors.pf007
\begin{proof}
	Set $\mathcal{I}_i:=\cate{K}^+(\mathcal{A}_i^\flat)$, $i=1,2$. The only
	point to prove is that, for every
	$X_1\in\Obj(\cate{K}^+(\mathcal{A}_1^\flat))$, the functor
	$(\cate{K}F)(X_1,\cdot):\cate{K}^+(\mathcal{A}_2^\flat)
	\to\cate{K}^+(\mathcal{B})$ sends acyclic complexes to acyclic complexes;
	the same statement holds after interchanging the indices $1$ and $2$. The
	existence of $\mathrm{R}F$ (respectively, $\mathrm{L}F$) then follows
	directly from Proposition~\ref{prop:localization-injective-bifunctor}.

	Take $Y\in\Obj(\cate{C}^+(\mathcal{A}_2^\flat))$. Then
	$(\cate{K}^2F)(X_1,Y)$ belongs to
	$\Obj(\cate{C}^2_f(\mathcal{B}))$. By condition (I2), each column
	$F(X_1^p,Y^\bullet)$ is acyclic. Therefore
	Corollary~\ref{prop:acyclic-assembly} implies that
	$\cate{K}F(X_1,Y):=\tot(\cate{K}^2F(X_1,Y))$ is exact, as desired.
\end{proof}

% segment-id: o014.li2.chapter4.bounded-derived-functors.p012
If $F$-injective (respectively, $F$-projective) subcategories are available,
the functor $\mathrm{R}F$ (respectively, $\mathrm{L}F$) can be determined by
\eqref{eqn:derived-bifunctor-1} (respectively,
\eqref{eqn:derived-bifunctor-2}).
